%% kapa: exam=1 round=24 year=2013 period=1 form=A subject=economics subject_ko=경제원론 seq=23 number=63
%% source: Q-Net 감정평가사 기출문제 (cst00309) artlSeq=5073802 file=extracted/1008/5073802/2013년 제24회 감정평가사 1차시험 기출문제/1교시(A형).hwp
%% conversion: kordoc 4.24.51 -> tools/gen1.py (auto); review=proofread-sonnet; flags=none
\begin{question}{63}
폐쇄경제의 IS-LM모형에서 지급준비율과 현금/예금 보유비율이 이자율의 감소함수일 때, 두 비율이 상수인 경우와 비교하여 옳은 설명을 모두 고른 것은? (단, IS곡선은 우하향, LM곡선은 우상향한다.)
\begin{framedtext}
ㄱ. 통화공급은 외생적으로 결정된다. \par
ㄴ. 통화정책의 효과가 커진다. \par
ㄷ. 재정정책의 효과가 커진다. \par
ㄹ. LM곡선의 기울기가 완만해진다.
\end{framedtext}
\begin{choices}
  \item ㄱ, ㄴ
  \item ㄴ, ㄷ
  \item ㄷ, ㄹ
  \item ㄱ, ㄴ, ㄹ
  \item ㄱ, ㄷ, ㄹ
\end{choices}
\end{question}
\answer{3}
